% =====================================================================
\section{Reflection and the Post Chain}\label{sec:reflection}
\pipestrip{8}
A rasteriser has no ray tracing, so reflection is faked from what is already on screen; no image or cube map is used.

\subsection{Specular ambient and wet surfaces}
\autofig{diagrams}{D2_fresnel_reflectivity}
Every shaded point adds a specular \emph{ambient} term $t_{sa}$ from the sky colour (modulated by ambient occlusion), so smooth metal shows a faint sky reflection even in the dark. Rain adds a water film on up-facing surfaces:
$\text{film}=w\cdot\mathrm{clamp}(1.4\,N_y+0.1,0,1)$ with the wetness $w\in[0,1]$; wet porous materials (concrete, wood, soil, HESCO, asphalt) darken by up to 45\,\%, roughness drops towards 0.3, and the ground shader mirrors the practical lights with a
per-fragment GGX highlight. The reflectivity stored for the next pass is $\max(0.45\,\text{film},\;m\,(1-\rho)^2\,0.6)$: wet ground and polished metal reflect, rough paint does not.

\subsection{Screen-space reflection (SSR)}
The main pass writes a second render target with the world normal and that reflectivity. A later pass (Listing~\ref{lst:Ssr}, key \texttt{F7}) reconstructs the world position $\vect P$ of each reflective pixel from depth, reflects the view direction, $\vect R=\mathrm{reflect}(\vect V,\vect N)$, and \emph{marches} $\vect Q(t)=\vect P+t\vect R$:
\begin{enumerate}
  \item \textbf{Ray march.} Start at $t=0.04(1+\text{jitter})$ (per-pixel jitter that TAA averages out) and grow the step, $t\leftarrow1.09\,t+0.03$, for 48 steps (64 on EPIC). Each sample is projected to the screen, $\vect q=\tfrac12\,\text{clip}_{xy}/\text{clip}_w+\tfrac12$;
  leaving the screen or going behind the camera ends the march.
  \item \textbf{Hit test.} The sample's linear depth (the clip $w$) is compared with the depth buffer at $\vect q$; a hit is a ray behind the surface by less than $0.15+0.04\,t+0.01\,d_0$ ($d_0$ = distance of $\vect P$ from the camera), growing with distance so far gaps are not missed.
  \item \textbf{Bisection.} Six bisection steps between the last miss $t_a$ and the hit $t_b$ find the crossing and remove stair-stepping.
  \item \textbf{Weights.} The fetched colour is multiplied by an \emph{edge fade} (smoothstep over 7\,\% of each border), a \emph{distance fade}, and the Fresnel factor
  $F=0.02+0.98(1-\vect N\!\cdot\!(-\vect V))^5$ mixed as $\mathrm{mix}(0.35,1,F)$.
\end{enumerate}
\lstSsr
\begin{figure}[h]
\centering
\piclab{ssr_on_off}{\linewidth}
\caption{Screen-space reflection off (left) and on (right) on wet ground: lamp posts and vehicles appear in the puddles (in-program render, key \texttt{F7}).}
\end{figure}
\begin{figure}[h]
\centering
\begin{subfigure}{0.48\linewidth}\picdia{ssr_raymarch}{\linewidth}\caption{ray march and bisection}\end{subfigure}
\begin{subfigure}{0.48\linewidth}\picdia{fresnel}{\linewidth}\caption{Schlick Fresnel vs.\ angle}\end{subfigure}
\caption{The geometry of the SSR march and the Fresnel weight.}
\end{figure}
Limits: SSR reflects only what is on screen, and reflections on rough surfaces are not blurred; there is no environment probe.

\subsection{The post chain in brief}
The linear HDR scene then passes through the chain of Chapter~\ref{sec:pipeline}: TAA, bloom, auto-exposure and \textbf{tone mapping} with the fitted ACES curve of Listing~\ref{lst:Aces}, $\mathrm{ACES}(x)=\frac{x(2.51x+0.03)}{x(2.43x+0.59)+0.14}$, followed by the sRGB transfer, vignette and grain.
\lstAces
\begin{figure}[h]
\centering
\piclab{pipeline_buffers}{\linewidth}
\caption{Intermediate buffers of one frame (in-program render): depth, normal, ambient occlusion, fog, HDR scene, bloom, final.}
\end{figure}
